\chapter{Introduction}
\textit{The true is the whole. However, the whole is only the essence completing itself through its own development.}\par
\begin{raggedleft}
    G.W.F Hegel: Phenomenology of Spirit \cite[p.11]{hegel1977pheno}\par
\end{raggedleft} \vspace{0.8cm}
The muon was not welcomed quite enthusiastically by particle physicists when it made its first appearance in cloud chamber experiments in the 1930s \cite{neddermeyer1937note}: \enquote{Who ordered that?} was Isidor Rabi's famous reaction to the strange new particle \cite{bartusiak1987who}, encapturing the disorientation of a scientific community who firmly believed matter to consist of electrons, neutrons and protons. The muon fitted nowhere in this scheme: Sharing all its properties with the electron except having a mass 200 times higher, the muon seemed to be basically a heavy copy of the electron. The mysterious sibling raised more questions than it answered.

But soon the muon took its place at the very heart of particle research, acting as a messenger for a plethora of unknown physics. Being only short-lived itself and quickly decaying into an electron, it is ecactly the spectrum of latter that provides the main subject for examination. As the produced electron varies both in energy and momentum at each decay instance, it follows that two more particles are being produced which carry away the missing quantities \cite{steinberger1949range}. Those particles were soon identified as neutrinos and in the Brookhaven experiment \cite{danby1962neutrinos} shown to exhibit the same sibling-like relation as the muon and the electron -- thus introducing the fundamental idea of the Standard Model that all matter particles exist as generations. But the electron spectrum becomes even more insightful with the parametrization that Louis Michel introduced \cite{michel1950interaction}, which turned it into a systematic probe of the underlying interaction. This mechanism for muon decay shares many characteristic with beta decay. When both processes were experimentally proven to maximally violate parity, the long-held premise that nature was mirror-symmetric, it was clear that they represent instances of the same theory: The universal vector-axial interaction \cite{sudarshan1958chirality,feynman1958theory} of weak theory. So muon decay arguably has served for many decades as almost a kind of Rosetta stone for particle physics. But -- has its translative potential truly been exhausted yet?

The apparently local interaction governing muon decay was eventually understood as the low-energy limit of the mediation by a $W$ boson. Invisible at the energy scale of muon decay, this heavy field becomes a dynamic degree of freedom only in the high-energy regime such as present in particle colliders \cite{arnison1983w,banner1983w} -- but leaves an imprint on the low-energy regime nevertheless: Its mass determines the decay rate of the muon. Could there be other heavy particles veiled by a similar fashion? In order to bridge the energy regimes and systematically probe for new physics, we perturb the standard set of Michel parameters $\rho,\eta,\eta'',\delta,\xi,\xi',\xi'',\tfrac{\alpha'}{A},\tfrac{\beta'}{A}$ and the Fermi constant $G_F$ with Low Energy Effective Field Theory (LEFT) operators. They represent the adequate physical description for interactions below the electroweak scale and by matching them onto the Standard Model Effective Field Theory (SMEFT) we trace the perturbations upward the energy regime. At the scale of new physics we open up the SMEFT operators and ask: Which heavy new fields could produce them? This approach allows us to systematically collect all Beyond Standard Model (BSM) fields which at tree-level produce effective operators that could significantly perturb muon decay. We constraint them by translating the experimental values of the Michel parameters into boundaries for the masses of the heavy new fields. In this way we hold yet another careful reading of the muonic Rosetta stone and ask: Has the illustrious sibling already revealed each of its secrets?

The structure of the thesis is as follows: The second chapter revolves around the Michel parameters of muon decay, whereby we first briefly introduce the Standard Model and its effective description of muon decay. Subsequently the parameters themselves are introduced and how they shed light on the underlying interaction of muon decay. We end the chapter with a discussion of the observational values for the Michel parameters and how to appropriately account for their correlations. 

In Chapter three the LEFT is introduced and we show how the presence of LEFT operators perturbs the Michel parameters. We present our calculations for the case of SM-like neutrino output and the full four-fermion results. This allows us to present constraints for various LEFT operators of dimension six. Finally the Renormalization group equations are applied to evolve the LEFT operators up to the electroweak scale.

As we reach the electroweak scale in chapter four, we are ready to delve into the Standard Model Effective Field Theory (SMEFT). After a brief introduction we match our LEFT operators of dimension six to SMEFT operators of dimensions six and seven. By using the Renormalization Group Equations we translate the LEFT boundaries to SMEFT boundaries both at the EW scale and at the scale of new physics. 

At the energy scale of new physics in chapter five we use the complete tree-level dictionary for BSM fields to SMEFT operators up to dimension seven to collect the possible UV origins for our set of SMEFT operators. 

In chapter six we present our results: The constraints for a charged electroweak scalar doublet and a correlated neutral vector mediator which produces multiple LEFT operators that perturb muon decay are derived. We discuss the effects an extension of the SM by each of the heavy particles would have and compare our constraints to the literature.

The outlook in chapter seven gives us the opportunity to reflect on the directive of our investigation. Starting from the low-energy regime and the experimental data for muon decay, our approach is effectively bottom-up: We expand the theoretical space of muon decay and trace the perturbations up to the ultra-violet -- and subsequently back down to muon decay. This contradicts the implicit assumption of many studies, starting from a concrete UV model, that the \enquote{true} theory is the one valid at all energies -- that the \textit{true} is the \textit{whole}, so to say. But experimental data and theoretical descriptions stretch a far more complex space than such a top-down system. For the course of this thesis we consequently focus on the translations between energy regimes and operator bases -- the renormalization group equations and matching conditions -- rather than on the UV models themselves. 
\section{Notation}
We use the following notation:
\begin{table}[h]
    \centering
    \begin{tabular}{ll}
         $L^{(d)}_i$ &  Wilson coefficients of the LEFT\\
         $C^{(d)}_i$ &  Wilson coefficients of the SMEFT\\
         $\mathcal{O}_X$ & Operators of the LEFT or the SMEFT\\
         $l$ & Lepton doublet\\
         $v$ & VEV\\
         $S_i$ & UV scalar\\
         $F_i$ & UV fermion\\
         $V_i$ & UV vector \\
         $q$ & Muon momentum\\
         $m_\mu$ & Muon mass\\
         $p$ & Electron momentum\\
         $E_e$ & Electron energy\\
         $m_e$ & Electron mass\\
         $k_i$ & Neutrino momenta
    \end{tabular}
    \caption{The notation used in this thesis.}
    \label{tab:notation}
\end{table}
